% Part: lambda-calculus
% Chapter: introduction
% Section: syntax

\documentclass[../../../include/open-logic-section]{subfiles}

\begin{document}

\olfileid[hi]{lam}{int}{syn}
\olsection{लैम्ब्डा कलन का वाक्यविन्यास}

हम चरों के अनुक्रम $x$, $y$, $z$,~\dots और कुछ अचर चिह्नों $a$, $b$,
$c$,~\dots से आरंभ करते हैं। पदों का समुच्चय निम्न प्रकार आगमनतः परिभाषित
होता है:
\begin{enumerate}
\item प्रत्येक चर एक पद है।
\item प्रत्येक अचर एक पद है।
\item यदि $M$ और $N$ पद हैं, तो $(MN)$ भी पद है।
\item यदि $M$ कोई पद और $x$ कोई चर है, तो $(\lambd[x][M])$ पद है।
\end{enumerate}
$(MN)$ रूप के पदों को \emph{अनुप्रयोग} और $(\lambd[x][M])$ रूप के पदों को
\emph{अमूर्तन} कहते हैं।

जिस पद्धति में कोई अचर नहीं होता उसे \emph{शुद्ध} लैम्ब्डा कलन कहते हैं।
हम मुख्यतः शुद्ध $\lambd$-कलन में काम करेंगे, इसलिए सभी छोटे लैटिन अक्षर
चरों को दर्शाएँगे। $\lambd$-कलन के पदों को दर्शाने के लिए हम बड़े लैटिन
अक्षरों ($M$, $N$, आदि) का उपयोग करते हैं।

हम संकेत-लिपि के कुछ नियम अपनाएँगे:

\begin{conv}
\begin{enumerate}
\item जब कोष्ठक छोड़ दिए जाते हैं, तब अनुप्रयोग बाएँ से दाएँ होता है।
  उदाहरणार्थ, यदि $M$, $N$, $P$ और $Q$ पद हैं, तो $MNPQ$, ~$(((MN)P)Q)$
  का संक्षिप्त रूप है।
\item फिर, जब कोष्ठक छोड़ दिए जाते हैं, तब लैम्ब्डा अमूर्तन को संभव
  अधिकतम व्याप्ति दी जाती है। उदाहरणार्थ, $\lambd[x][MNP]$ को
  $(\lambd[x][((MN)P)])$ पढ़ा जाता है।
\item एक लैम्ब्डा से अनेक चरों का अमूर्तन किया जा सकता है। उदाहरणार्थ,
  $\lambd[xyz][M]$, ~$\lambd[x][\lambd[y][\lambd[z][M]]]$ का संक्षिप्त रूप
  है।
\end{enumerate}
\end{conv}

उदाहरणार्थ,
\[
\lambd[xy][xxyx \lambd[z][xz]]
\]
इसका संक्षिप्त रूप है:
\[
\lambd[x][\lambd[y][((((xx)y)x)(\lambd[z][(xz)]))]].
\]
आपको ये नियम याद कर लेने चाहिए। आरंभ में ये आपको चकरा देंगे, किंतु आप इनके
अभ्यस्त हो जाएँगे; और कुछ समय बाद ये कोष्ठकों के जंगल से जूझने की अपेक्षा
कम परेशान करेंगे।

जो दो पद केवल अपने आबद्ध चरों के नामों में भिन्न हों, उन्हें
$\alpha$-तुल्य कहते हैं; उदाहरणार्थ, $\lambd[x][x]$ और $\lambd[y][y]$।
इनको ``एक ही'' पद मानना सुविधाजनक होगा; दूसरे शब्दों में, जब हम कहते हैं
कि $M$ और $N$ समान हैं, तो हमारा अर्थ ``आबद्ध चरों के पुनर्नामकरण तक'' भी
होता है। जो चर किसी $\lambd$ की व्याप्ति में हों वे ``आबद्ध'' और अन्य चर
``मुक्त'' कहलाते हैं। पिछले उदाहरण में कोई मुक्त चर नहीं है; किंतु
\[
(\lambd[z][yz])x
\]
में $y$ और $x$ मुक्त हैं तथा $z$ आबद्ध है।

\end{document}
